\section{CONCLUSION}

Bài báo này đề xuất phương pháp mạng từ trở sinh lưới trong không gian ba chiều (3D-MBGRN) dựa trên hệ tọa độ trụ cho bài toán phân tích từ trường phi tuyến của máy điện quay đối xứng, với mô hình kiểm chứng là động cơ nam châm vĩnh cửu hướng trục (AFPM) 30 rãnh / 20 cực. Bằng việc rời rạc hóa miền tính toán thành tập hợp các phần tử lục diện đồng nhất và áp dụng kỹ thuật dịch chỉ số lưới (index-shifting), phương pháp mô hình hóa chuyển động quay của rotor trên cấu trúc lưới cố định mà không đòi hỏi các tác vụ tái tạo lưới hay nội suy phức tạp tại mặt trượt. Bên cạnh đó, bộ giải phi tuyến lặp điểm cố định tích hợp thuật toán Adaptive Material Decay kết hợp cơ chế quay lui (backtracking) đã giải quyết triệt để hiện tượng bẫy trì trệ và dao động số học, đảm bảo quá trình hội tụ đơn điệu và bền vững.

Ý nghĩa khoa học của phương pháp được khẳng định thông qua việc đối chiếu trực tiếp với mô hình 3D-FEM ở cả trạng thái không tải và có tải định mức. Dù vận hành trên hai nền tảng giải số và kỹ thuật hậu xử lý hoàn toàn độc lập, tất cả các đại lượng điện từ cốt lõi ($B_z, \Psi_a, e_a, T, F_z$) đều đạt sự tương đồng cao về dạng sóng transient và phổ hài, với độ lệch duy trì trong khoảng từ $0.3\%$ đến $3.1\%$. Đáng chú ý, cấu trúc lưới có tổ chức của 3D-MBGRN giúp triệt tiêu các dao động số học cục bộ vốn thường xuất hiện trong 3D-FEM do ảnh hưởng của tái chia lưới phi cấu trúc.

Về mặt hiệu năng tính toán, phân tích hội tụ số trị chứng minh bộ giải 3D-MBGRN có khả năng đạt độ chính xác tương đương với mô hình phần tử hữu hạn phân giải cao trong khi giảm thiểu đáng kể chi phí tài nguyên. Nhờ cấu trúc ma trận thưa 7-point stencil và việc gán ẩn số thế từ duy nhất tại tâm ô lưới lục diện, bậc tự do (DoF) của hệ phương trình giảm khoảng $1.4\text{ lần}$ so với 3D-FEM ở cùng quy mô số lượng phần tử. Đặc tính này kết hợp với kỹ thuật dịch chỉ số giúp bộ giải tiết kiệm dung lượng bộ nhớ RAM và rút ngắn thời gian mô phỏng transient từ $6$ đến hơn $16\text{ lần}$. Hiệu quả tính toán vượt trội này khẳng định 3D-MBGRN là một công cụ phân tích từ trường tiềm năng, có khả năng bổ trợ hoặc thay thế hiệu quả cho 3D-FEM trong các bài toán thiết kế và tối ưu hóa đa mục tiêu quy mô lớn.

Trong các nghiên cứu tiếp theo, phương pháp 3D-MBGRN có thể tiếp tục được phát triển và hoàn thiện theo ba định hướng trọng tâm:
\begin{itemize}
    \item \textbf{Mở rộng áp dụng cho các cấu hình động cơ khác nhau:} Áp dụng và kiểm chứng mô hình trên nhiều dòng máy điện có cấu trúc hình học và đường đi từ thông 3D phức tạp, chẳng hạn như động cơ nam châm chìm (IPMSM), động cơ từ trở đồng bộ (SynRM), máy điện dòng từ hỗn hợp (Transverse Flux Machine) hoặc các cấu trúc máy điện hướng trục có rãnh/nam châm chéo góc (skewing).
    \item \textbf{Tích hợp mô hình tính toán tổn hao:} Phát triển các thuật toán hậu xử lý dựa trên trường từ thông cục bộ $B(t)$ để trích xuất chính xác tổn hao sắt từ (tổn hao trễ và dòng xoáy trong lõi thép gông/răng stator theo các mô hình Bertotti hay Steinmetz), tổn hao dòng xoáy bên trong khối nam châm vĩnh cửu và tổn hao đồng cục bộ do hiệu ứng bề mặt và hiệu ứng lân cận (skin and proximity effects).
    \item \textbf{Tối ưu hóa và phân tích đa vật lý (Multiphysics Optimization):} Tận dụng trực tiếp dữ liệu nguồn nhiệt từ tổn hao và lực điện từ trích xuất từ bộ giải 3D-MBGRN để ghép nối đồng thời với mô hình mạng nhiệt tương đương (LPTN)/FEM nhiệt và mô hình phân tích dao động - tiếng ồn (NVH), từ đó tích hợp vào các thuật toán tối ưu hóa đa mục tiêu phục vụ quy trình thiết kế động cơ toàn diện.
\end{itemize}